% A test document for the paper pages (papers/README.md). Not a real paper:
% it exercises what papers usually contain, plus one of each construct that
% prepare.lua and paper.cfg exist for (marked "test:"), so changes there, to the
% template or to the CSS can be checked against it.
\documentclass[11pt]{article}
\usepackage{amsmath,amssymb,amsthm}
\usepackage{natbib}
\usepackage{hyperref}

\newtheorem{definition}{Definition}
\newtheorem{proposition}{Proposition}
\newcommand{\Alg}{\mathcal{A}}
\newcommand{\Hilb}{\mathcal{H}}
% test: physics-style macros, as in older papers
\newcommand{\tn}[1]{\mbox{\tiny #1}}
\newcommand{\GN}{G_{\tn N}}       % a _ inside a preamble macro
\newcommand{\cB}{{\cal B}}        % old-style font switch
\newcommand{\beq}{\begin{eqnarray}}  % environment shorthands
\newcommand{\eeq}{\end{eqnarray}}

\title{Rendering Test: A Paper as a Web Page}
\author{Enrico Cinti}
\date{}

\begin{document}
\maketitle

\begin{abstract}
This document is a fixture for the website's paper pages. It contains the
pieces a philosophy of physics paper tends to have (numbered sections,
inline and display mathematics, a definition and a proposition, footnotes,
citations and a reference list) so that the conversion from LaTeX to HTML
can be checked in one place.
\end{abstract}

\section{Introduction}

Dualities relate theories that look very different but agree on all their
physical content \citep{deharo2016}. A standard example is position--momentum
duality in quantum mechanics, where the Fourier transform
$\hat{f}(p) = \int e^{-ipx} f(x)\,dx$ maps one representation onto the
other.\footnote{Footnotes are collected at the end of the page and linked
both ways.} Whether such theories are equivalent is a philosophical question,
not only a formal one \citep{butterfield2021}.

\subsection{What the page should show}

Sections are numbered as in the PDF, and the table of contents on the right
follows the reader. Cross-references work too: see
Section~\ref{sec:algebras} and Equation~\eqref{eq:commutant}.

\section{Algebras of observables}
\label{sec:algebras}

Let $\Alg$ be a von Neumann algebra acting on a Hilbert space $\Hilb$. Its
commutant is
\begin{equation}
  \Alg' = \{\, b \in B(\Hilb) : ab = ba \ \text{for all } a \in \Alg \,\},
  \label{eq:commutant}
\end{equation}
and the bicommutant theorem says that $\Alg'' = \Alg$.

\begin{definition}
A state $\omega$ on $\Alg$ is \emph{faithful} if $\omega(a^*a) = 0$
implies $a = 0$.
\end{definition}

\begin{proposition}
If $\Omega \in \Hilb$ is cyclic and separating for $\Alg$, the vector state
$\omega(a) = \langle \Omega, a\,\Omega \rangle$ is faithful.
\end{proposition}

Older notation comes through as well: Newton's constant $\GN$, an
algebra $\cB$, and an \texttt{eqnarray} written with shorthands, with an
accent and a superscript on bare macros,
\beq\label{eq:crossed}
\hat\Alg \;=\; \Alg \rtimes_{\sigma^\Omega} \mathbb{R}~,
\eeq
referred to as Equation~\eqref{eq:crossed}.

Aligned equations render as well:
\begin{align}
  S &= \int d^4x \sqrt{-g}\,\Big( \frac{R - 2\Lambda}{16\pi G} + \mathcal{L}_m \Big), \\
  G_{\mu\nu} + \Lambda g_{\mu\nu} &= 8\pi G\, T_{\mu\nu}.
\end{align}

\section{Conclusion}

If this page reads well, so will the real ones.\footnote{And if it does not,
the fix goes in \texttt{render.sh} or \texttt{paper.html}, not in the paper.}

\bibliographystyle{plainnat}
\bibliography{refs}
\end{document}
